\chapter[İçerme-Dışarma · \textit{Orta}]{İçerme-Dışarma}
\chaptermark{İçerme-Dışarma}

Bölüm 3'te ele aldığımız temel sayma kurallarından ilki toplama
kuralıydı: Eğer iki küme birbirinden bütünüyle ayrıksa, bu iki kümenin
birleşiminin eleman sayısı, kümelerin eleman sayıları toplamına eşittir. Ancak
uygulamada ve algoritmik problemlerde karşımıza çıkan kümeler nadiren bu denli
birbirinden yalıtılmıştır. Kümeler sıklıkla kesişir, ortak elemanlar barındırır
ve bu ortak elemanlar üst üste biner.

Birbirini kesen kümelerin birleşimini sayarken doğrudan toplama yapmak
bizi hataya sürükler; çünkü ortak elemanları birden fazla kez saymış oluruz.
Bu bölümde, üst üste binen kümelerin eleman sayısını hatasız bir şekilde
belirlememizi sağlayan temel yöntemlerden birini, \textbf{içerme-dışarma}
(inclusion--exclusion) ilkesini inşa edeceğiz.

\section{İki ve Üç Küme}

\subsection{İki Küme ve Çift Sayımı Düzeltme}

İlkenin mantığını somut bir soru üzerinden görelim.

\begin{example} 1'den 100'e kadar olan tam sayılar arasından 2 veya 3 ile
bölünebilen kaç sayı vardır? \end{example}

\begin{proof}[Çözüm]
2'ye bölünen tam sayıların kümesine $A$, 3'e bölünen tam sayıların kümesine
de $B$ diyelim. Amacımız bu iki kümenin birleşimi olan $A \cup B$ kümesinin eleman
sayısını, yani $|A \cup B|$ değerini bulmaktır.

1 ile 100 arasında 2'ye bölünen tam sayılar $2, 4, 6, \dots, 100$ biçimindedir. Bu
sayıların adedi: \[ |A| = \left\lfloor \frac{100}{2} \right\rfloor = 50 \] Benzer
biçimde, 3'e bölünen tam sayılar $3, 6, 9, \dots, 99$ biçimindedir ve adedi: \[ |B|
= \left\lfloor \frac{100}{3} \right\rfloor = 33 \] Bu iki değeri doğrudan toplarsak
$50 + 33 = 83$ buluruz; ancak bu toplam doğru sonucu vermez.

Küçük bir sayıyı, örneğin 6'yı ele alalım. 6 sayısı çift olduğu için $A$ kümesinin
içindedir; dolayısıyla 50'nin içinde bir kez sayılmıştır. Öte yandan 6 sayısı
3'ün de katı olduğundan $B$ kümesinin içindedir ve 33'ün içinde de bir kez
sayılmıştır. Demek ki 6 sayısı bu 83 toplamına bir değil, iki kez dahil
edilmiştir. Yalnızca 6 değil; $12, 18, 24, \dots$ gibi hem 2'ye hem de 3'e bölünen
tüm sayılar ikişer kez sayılmıştır.

2 ve 3 aralarında asal olduğundan (Bölüm 11 ve 12'deki bölünebilme ilkeleri
uyarınca), bir sayının hem 2'ye hem de 3'e bölünebilmesi o sayının 6'ya tam
bölünmesi demektir. Bu ortak sayıların kümesi tam olarak $A \cap B$ kesişimidir.
1 ile 100 arasında 6'ya bölünen tam sayıların adedi: \[ |A \cap B| = \left\lfloor
\frac{100}{6} \right\rfloor = 16 \] Her bir ortak elemanı fazladan birer kez
saydığımıza göre, bu fazlalığı gidermenin yolu kesişimdeki eleman sayısını
toplamdan bir kez çıkarmaktır. O hâlde aranan yanıt: \[ |A \cup B| = |A| + |B| -
|A \cap B| = 50 + 33 - 16 = 67 \] olarak elde edilir.
\end{proof}

Çözümü bir Venn şemasıyla özetleyelim; içerme-dışarmanın neyi düzelttiği
şekilde tek bakışta görülür:

\begin{figure}[h]
\centering
\begin{tikzpicture}[scale=0.9]
  \draw[kitap-cizgi, fill=black!3] (-2.6,-1.9) rectangle (2.6,1.9);
  \node[etiket, anchor=north west] at (-2.5,1.83) {\itshape E};
  \begin{scope}
    \clip (0.65,0) circle (1.35);
    \fill[kitapvurguacik] (-0.65,0) circle (1.35);
  \end{scope}
  \draw[kitap-cizgi] (0.65,0) circle (1.35);
  \draw[kitap-cizgi] (-0.65,0) circle (1.35);
  \node[etiket] at (-0.72,1.62) {$A$};
  \node[etiket] at (0.72,1.62) {$B$};
  \node[etiket] at (0,0.05) {16};
  \node[etiket] at (-1.3,-0.2) {34};
  \node[etiket] at (1.3,-0.2) {17};
  \node[etiket] at (0,-1.55) {33};
\end{tikzpicture}
\caption{İki küme için içerme-dışarma: kesişimdeki 16 sayı, $50 + 33$ toplamında
iki kez sayıldığı için bir kez çıkarılır; $|A \cup B| = 50 + 33 - 16 = 67$.
Dışarıdaki 33 sayı ne 2'ye ne de 3'e bölünür.}
\end{figure}

Bu basit örnek genel kuralın özünü gösterir: Bir elemanı iki kez saydıysak,
net sayımını 1'e indirmek için onu bir kez dışarı çıkarmalıyız (dışarma).

\begin{theorem}[İki Küme İçin İçerme-Dışarma İlkesi] $A$ ve $B$ sonlu iki küme olmak
üzere, \[ |A \cup B| = |A| + |B| - |A \cap B| \] eşitliği geçerlidir.
\end{theorem}

\begin{proof} $A \cup B$ kümesindeki herhangi bir $x$ elemanını ele alalım.
Bu eleman için tam olarak üç durum söz konusudur:
\begin{enumerate} \item $x \in A$ fakat $x \notin B$ (yalnızca $A$'da olanlar): Bu
durumda $x$ elemanı $|A|$ içinde 1 kez, $|B|$ içinde 0 kez, $|A \cap B|$ içinde 0 kez
sayılır. Eşitliğin sağ tarafındaki net katkısı $1 + 0 - 0 = 1$'dir. \item $x \in B$
fakat $x \notin A$ (yalnızca $B$'de olanlar): Bu durumda $x$ elemanı $|A|$ içinde 0 kez,
$|B|$ içinde 1 kez, $|A \cap B|$ içinde 0 kez sayılır. Eşitliğin sağ tarafındaki net
katkısı $0 + 1 - 0 = 1$'dir. \item $x \in A$ ve $x \in B$ (ortak olanlar): Bu durumda
$x$ elemanı $|A|$ içinde 1 kez, $|B|$ içinde 1 kez ve $|A \cap B|$ içinde 1 kez sayılır.
Eşitliğin sağ tarafındaki net katkısı $1 + 1 - 1 = 1$'dir. \end{enumerate} Her
durumda birleşime ait olan herhangi bir eleman sağ taraftaki formülde net olarak
tam 1 kez sayılmaktadır. Dolayısıyla eşitlik her zaman doğrudur. \end{proof}

\subsection{Üç Kümeye Geçiş: Dengeyi Kurmak}

Şimdi üçüncü bir koşul ekleyerek adımları genişletelim.

\begin{example} 1'den 100'e kadar olan tam sayılar arasından 2, 3 veya 5 ile
bölünebilen kaç sayı vardır? \end{example}

\begin{proof}[Çözüm]
2'ye bölünenlerin kümesi $A$, 3'e bölünenlerin kümesi $B$ ve 5'e bölünenlerin kümesi
$C$ olsun. Amacımız $|A \cup B \cup C|$ değerini bulmaktır.

Önce her kümenin büyüklüğünü tek tek hesaplayalım: \begin{align*} |A| &=
\left\lfloor \frac{100}{2} \right\rfloor = 50 \\ |B| &= \left\lfloor
\frac{100}{3} \right\rfloor = 33 \\ |C| &= \left\lfloor \frac{100}{5}
\right\rfloor = 20 \end{align*} Tüm bu elemanları doğrudan toplarsak: \[ |A| +
|B| + |C| = 50 + 33 + 20 = 103 \] elde ederiz. 1'den 100'e kadar zaten 100 sayı
varken sonucun 103 çıkması, kimi elemanların birden fazla kez sayıldığını gösterir.

İkişerli ortak elemanları ele alalım: \begin{itemize} \item $2 \times 3 = 6$ ile
bölünenler hem $A$'da hem $B$'dedir: $|A \cap B| = \lfloor 100/6 \rfloor = 16$. \item
$2 \times 5 = 10$ ile bölünenler hem $A$'da hem $C$'dedir:
$|A \cap C| = \lfloor 100/10 \rfloor = 10$. \item $3 \times 5 = 15$ ile bölünenler
hem $B$'de hem $C$'dedir: $|B \cap C| = \lfloor 100/15 \rfloor = 6$. \end{itemize} İki
küme probleminde yaptığımız gibi, ikişerli kesişimleri dışarı atalım: \[ 103 -
(16 + 10 + 6) = 103 - 32 = 71 \] Peki işlem bu aşamada tamamlandı mı?

30 sayısına bakalım. 30, $2 \times 3 \times 5 = 30$ olduğundan hem 2'ye, hem
3'e, hem de 5'e tam bölünür; yani $30 \in A \cap B \cap C$ kümesindedir. 30'un
yukarıdaki işlemlerde kaç kez hesaba katıldığını kontrol edelim:
\begin{itemize} \item $|A|$, $|B|$ ve $|C|$ toplamında 3 kez eklendi (+3). \item
$|A \cap B|$, $|A \cap C|$ ve $|B \cap C|$ çıkarılırken 3 kez çıkarıldı (-3).
\end{itemize} Net durum: $+3 - 3 = 0$. Yani 30 sayısı (ve onun gibi 2, 3, 5'in ortak
katları olan 60 ve 90), elde ettiğimiz 71 sonucunun içinde \textbf{hiç
sayılmamıştır}. Birleşime ait bir elemanı bütünüyle dışarıda bırakmış olduk.

Bu eksikliği gidermek ve bu elemanları doğru biçimde saymak için, üç kümenin
kesişimini tekrar \textbf{içeri almalı}, yani eklemeliyiz: \[ |A \cap B \cap C| =
\left\lfloor \frac{100}{30} \right\rfloor = 3 \] Böylece doğru sayım: \[ |A \cup B
\cup C| = 71 + 3 = 74 \] olarak bulunur. \end{proof}

Bu sayımın dengesini Venn şeması üzerinde de okuyabiliriz; şekildeki her
bölgeye o bölgedeki sayı adedi yazılmıştır:

\begin{figure}[h]
\centering
\begin{tikzpicture}[scale=0.95]
  \draw[kitap-cizgi, fill=black!3] (-1.8,-1.7) rectangle (3.4,3.1);
  \node[etiket, anchor=north west] at (-1.7,3.03) {\itshape E};
  \begin{scope}
    \clip (0,0) circle (1.5);
    \clip (1.6,0) circle (1.5);
    \fill[kitapvurguacik] (0.8,1.4) circle (1.5);
  \end{scope}
  \draw[kitap-cizgi] (0,0) circle (1.5);
  \draw[kitap-cizgi] (1.6,0) circle (1.5);
  \draw[kitap-cizgi] (0.8,1.4) circle (1.5);
  \node[etiket] at (-1.05,1.9) {$A$};
  \node[etiket] at (2.65,1.9) {$B$};
  \node[etiket] at (0.8,2.72) {$C$};
  \node[etiket] at (0.8,0.55) {3};
  \node[etiket] at (0.8,-0.95) {13};
  \node[etiket] at (0.15,1.15) {7};
  \node[etiket] at (1.45,1.15) {3};
  \node[etiket] at (-1.0,0.25) {27};
  \node[etiket] at (2.6,0.25) {14};
  \node[etiket] at (0.8,2.32) {7};
  \node[etiket] at (3.15,1.0) {26};
\end{tikzpicture}
\caption{Üç küme için içerme-dışarma: bölgelerin toplamı
$27+14+7+13+3+7+3 = 74$'tür; aynı değer formülden
$50+33+20-(16+10+6)+3=74$ olarak gelir. Üçlü kesişimdeki 3 sayı, son adımda
geri eklenir.}
\end{figure}

Bu döngü içerme-dışarmanın temel mekanizmasıdır: Teklileri ekle (içerme),
çiftlileri çıkar (dışarma), üçlüleri geri ekle (içerme).

\begin{theorem}[Üç Küme İçin İçerme-Dışarma İlkesi] $A$, $B$ ve $C$ sonlu üç küme
olmak üzere, \[ |A \cup B \cup C| = |A| + |B| + |C| - (|A \cap B| + |A \cap C| +
|B \cap C|) + |A \cap B \cap C| \] eşitliği geçerlidir. \end{theorem}

\begin{proof} $A \cup B \cup C$ kümesindeki herhangi bir $x$ elemanının formülün sağ
tarafına yaptığı net katkıyı, $x$'in tam olarak kaç kümenin elemanı olduğuna göre
inceleyelim: \begin{enumerate} \item $x$ bu üç kümeden \textbf{tam birine} ait
olsun (örneğin sadece $A$'da): Tekli toplamlardan $|A|$'dan 1 katkı gelir. İkili
veya üçlü kesişimlerde yer almadığından diğer terimlerden 0 gelir. Net katkı: 1.
\item $x$ bu üç kümeden \textbf{tam ikisine} ait olsun (örneğin $A$ ve $B$'de, ama
$C$'de değil): Teklilerden $|A|$ ve $|B|$ olmak üzere 2 katkı gelir. İkililerden
yalnızca $|A \cap B|$ içinde yer aldığından -1 katkı gelir. Üçlü kesişimde yer
almaz (0). Net katkı: $2 - 1 = 1$. \item $x$ bu üç kümenin \textbf{üçüne de} ait
olsun ($A \cap B \cap C$): Tekli terimlerin üçünde de yer alır (+3). İkili
kesişimlerin üçünde de yer alır (-3). Üçlü kesişimde yer alır (+1). Net katkı:
$3 - 3 + 1 = 1$. \end{enumerate} Hangi durumda olursa olsun, birleşime ait her
eleman formül tarafından tam 1 kez sayılmaktadır. Dolayısıyla eşitlik
doğrudur. \end{proof}

\subsection{Çözümlü Örnekler}

\begin{example} 40 kişilik bir sınıfta 22 öğrenci matematik, 18 öğrenci fizik,
15 öğrenci kimya dersini almaktadır. 9 öğrenci hem matematik hem fizik, 7
öğrenci hem matematik hem kimya, 5 öğrenci ise hem fizik hem kimya dersini
almaktadır. Her üç dersi de alan 3 öğrenci olduğuna göre, bu üç dersten
\textbf{en az birini} alan kaç öğrenci vardır? Hiçbir dersi almayan kaç öğrenci
vardır? \end{example}

\begin{proof}[Çözüm]
Matematik alanların kümesi $M$, fizik alanların kümesi $F$, kimya alanların kümesi $K$
olsun. Verilen büyüklükler şunlardır: \begin{align*} |M| &= 22, \quad |F|
= 18, \quad |K| = 15 \\ |M \cap F| &= 9, \quad |M \cap K| = 7, \quad |F \cap K|
= 5 \\ |M \cap F \cap K| &= 3 \end{align*} En az bir dersi alan öğrenciler
$M \cup F \cup K$ kümesini oluşturur. Üç küme içerme-dışarma formülünü uygularız:
\begin{align*} |M \cup F \cup K| &= (|M| + |F| + |K|) \\ &\quad - (|M \cap F| + |M \cap K| + |F \cap K|) + |M \cap F \cap K| \\ &= (22 + 18 + 15) - (9 + 7 + 5) + 3 \\ &= 55 - 21 + 3 = 37 \end{align*} Demek ki öğrencilerin 37'si en az bir
ders almaktadır.

Sınıf mevcudu 40 olduğuna göre, hiçbir dersi almayan öğrenci sayısı evrensel
kümeden birleşimin çıkarılmasıyla bulunur: \[ 40 - |M \cup F \cup K| = 40 - 37
= 3 \] olarak elde edilir. \end{proof}

\begin{example} 1'den 600'e kadar olan tam sayılar arasından 4 veya 6 ile
bölünen fakat 5 ile \textbf{bölünmeyen} kaç sayı vardır? \end{example}

\begin{proof}[Çözüm]
İstenen koşul iki parçadan oluşur: Sayı 4 veya 6'ya bölünmeli, fakat 5'e
bölünmemelidir. Bu tip sorularda 4 veya 6'ya bölünenleri bulup içlerinden 5'e
bölünenleri çıkarmak en güvenli yoldur.

4 ile bölünenler $A$, 6 ile bölünenler $B$, 5 ile bölünenler $C$ olsun. Bizden istenen
küme $(A \cup B) \setminus C$ kümesidir. Kümeler teorisinden: \[ |(A \cup B)
\setminus C| = |A \cup B| - |(A \cup B) \cap C| \] yazabiliriz. Dağılma
özelliğinden $(A \cup B) \cap C = (A \cap C) \cup (B \cap C)$ olur.

Adım adım hesaplayalım: \begin{enumerate} \item Önce $A \cup B$ kümesini bulalım.
Hem 4'e hem 6'ya bölünen en küçük pozitif tam sayı (Bölüm 13'te ele aldığımız
EKOK tanımıyla) $\text{EKOK}(4, 6) = 12$'dir. \[ |A| = \left\lfloor \frac{600}{4}
\right\rfloor = 150, \quad |B| = \left\lfloor \frac{600}{6} \right\rfloor = 100,
\quad |A \cap B| = \left\lfloor \frac{600}{12} \right\rfloor = 50 \] Buradan: \[ |A
\cup B| = 150 + 100 - 50 = 200 \] \item Şimdi bu sayıların arasından 5'e de
bölünenleri, yani $(A \cap C) \cup (B \cap C)$ kümesini bulalım:
\begin{itemize} \item $A \cap C$: Hem 4'e hem 5'e bölünenler ($\text{EKOK}(4, 5) = 20$'ye
bölünenler): \[ |A \cap C| = \left\lfloor \frac{600}{20} \right\rfloor = 30 \]
\item $B \cap C$: Hem 6'ya hem 5'e bölünenler ($\text{EKOK}(6, 5) = 30$'a
bölünenler): \[ |B \cap C| = \left\lfloor \frac{600}{30} \right\rfloor = 20 \]
\item $(A \cap C) \cap (B \cap C) = A \cap B \cap C$: Hem 4, hem 6, hem 5'e bölünenler.
$\text{EKOK}(4, 6, 5) = 60$'a bölünür: \[ |A \cap B \cap C| = \left\lfloor
\frac{600}{60} \right\rfloor = 10 \] \end{itemize} İki küme formülünü $(A \cap C)$
ve $(B \cap C)$ için uygularsak: \[ |(A \cap C) \cup (B \cap C)| = 30 + 20 - 10 =
40 \] \item Son olarak istenen farkı alırız: \[ |(A \cup B) \setminus C| = 200 -
40 = 160 \] \end{enumerate} Demek ki 1'den 600'e kadar aradığımız özellikte 160
sayı vardır. \end{proof}

\begin{example} 6 basamaklı, rakamları toplamı 4 olan kaç farklı doğal sayı
yazılabilir? \end{example}

\begin{proof}[Çözüm]
Bir doğal sayının 6 basamaklı olması için ilk basamağının sıfırdan farklı olması
gerekir. Sayımız $a_1 a_2 a_3 a_4 a_5 a_6$ olsun. Buradaki kısıtlar şunlardır: \[
a_1 + a_2 + a_3 + a_4 + a_5 + a_6 = 4 \] ve $a_1 \ge 1$, diğer $a_i \ge 0$ olmalıdır.

İlk basamağın sıfır olmaması kuralını uygulamak için $a_1' = a_1 - 1$
tanımlayalım. Bu durumda $a_1' \ge 0$ olur ve denklemimiz: \[ a_1' + a_2 + a_3 +
a_4 + a_5 + a_6 = 4 - 1 = 3 \] hâline gelir. Burada her bir değişken negatif
olmayan birer tam sayıdır ($a_i \ge 0$). Normalde her bir basamak en fazla 9
değerini alabilir; ancak toplamları zaten 3 olduğu için hiçbir değişkenin 9'u
aşma tehlikesi yoktur. Dolayısıyla basamakların üst sınır kısıtı kendiliğinden
sağlanır.

Bölüm 5'te gördüğümüz ayraç (stars and bars) yöntemine göre, $n=3$ özdeş nesneyi
$k=6$ farklı kutuya dağıtma sayısı: \[ \binom{n + k - 1}{k - 1} = \binom{3 + 6
- 1}{6 - 1} = \binom{8}{5} = \binom{8}{3} = \frac{8 \times 7 \times 6}{3
\times 2 \times 1} = 56 \] olarak bulunur.

Üst sınır kısıtlarının toplamı aşmadığı durumlarda içerme-dışarmaya gerek
kalmadan doğrudan çözüme ulaşılır; ancak toplamın 9'u aştığı durumlarda
içerme-dışarma zorunlu hâle gelir. \end{proof}

\section{Genel Formül}

\subsection{Daha Fazla Küme ve Değişen İşaretler Deseni}

İki kümede tek bir çıkarma yaptık: $+ -$. Üç kümede bir çıkarma, bir toplama
yaptık: $+ - +$. Dört, beş veya genel olarak $n$ küme olduğunda da benzer bir yol
izleriz.

Her adımda birleşim alanını dengelemek için bir sonraki seviyedeki tüm kesişimleri
hesaba katmamız gerekir. Bu dalgalı işaret deseni (önce ekle, sonra çıkar, tekrar
ekle...) kombinatorik saymanın temel yapılarından biridir.

\begin{theorem}[Genel İçerme-Dışarma İlkesi] $A_1, A_2, \dots, A_n$ sonlu kümeler
olsun. Bu kümelerin birleşiminin eleman sayısı:
\begin{multline*}
|A_1 \cup A_2 \cup \cdots \cup A_n| = \sum_{1 \le i \le n} |A_i| - \sum_{1 \le i < j \le n} |A_i \cap A_j| \\
+ \sum_{1 \le i < j < k \le n} |A_i \cap A_j \cap A_k| - \cdots + (-1)^{n+1} |A_1 \cap A_2 \cap \cdots \cap A_n|
\end{multline*}
toplamı ile verilir. Toplam sembolüyle ifade edecek olursak: \[ \left|
\bigcup_{i=1}^n A_i \right| = \sum_{k=1}^n (-1)^{k+1} \sum_{1 \le i_1 < i_2 <
\dots < i_k \le n} |A_{i_1} \cap A_{i_2} \cap \dots \cap A_{i_k}| \] olur.
\end{theorem}

\subsection{Bu Formül Neden Her Zaman Doğrudur?}

Formülün geçerliliğini görmek için birleşime ait \textbf{tek bir elemanın}
adımlardaki durumunu izlemek yeterlidir.

Birleşim kümesinden herhangi bir $x \in \bigcup_{i=1}^n A_i$ elemanı alalım. Bu
eleman, verilen $n$ kümeden \textbf{tam $m$ tanesinin} içinde bulunsun
($1 \le m \le n$).

Formülün sağ tarafındaki terimlerin $x$ elemanını kaç kez saydığını inceleyelim:
\begin{enumerate} \item \textbf{Tekli kümeler ($\sum |A_i|$):} $x$ tam $m$ kümenin
içinde olduğuna göre, bu toplama tam $m = \binom{m}{1}$ kez katkı sağlar
($+ \binom{m}{1}$). \item \textbf{İkili kesişimler ($\sum |A_i \cap A_j|$):}
$x$'in bir $A_i \cap A_j$ kesişiminde yer alması için, hem $A_i$'nin hem de
$A_j$'nin $x$'i içeren o $m$ küme arasından seçilmiş olması gerekir. $m$ küme
arasından 2 küme $\binom{m}{2}$ farklı şekilde seçilebilir. Dolayısıyla bu
terimden $-\binom{m}{2}$ katkı gelir. \item \textbf{Üçlü kesişimler
($\sum |A_i \cap A_j \cap A_k|$):} Benzer şekilde, $x$ elemanı bu $m$ küme
arasından seçilen herhangi üçlünün kesişimindedir. Bu terimden $+\binom{m}{3}$
katkı gelir. \item Genel olarak, $k$'lı kesişimlerin toplamından $x$'e gelen
katkı $(-1)^{k+1}\binom{m}{k}$ olacaktır. Eğer $k > m$ ise, $x$ en fazla $m$
kümede yer aldığından daha büyük kesişimlerde bulunamaz; bu terimlerden gelen
katkı 0'dır. \end{enumerate}

O hâlde formülün sağ tarafının $x$ elemanına verdiği \textbf{net sayım değeri}:
\[ \text{Net Sayım} = \binom{m}{1} - \binom{m}{2} + \binom{m}{3} - \binom{m}{4} +
\cdots + (-1)^{m+1}\binom{m}{m} \] olur. Bu toplamın değerini bulmak için
Bölüm 6'da gördüğümüz binom açılımını hatırlayalım: \[ (1 - 1)^m = \binom{m}{0}
- \binom{m}{1} + \binom{m}{2} - \binom{m}{3} + \cdots + (-1)^m \binom{m}{m} \] Sol
taraf $(0)^m = 0$ olduğundan: \[ 0 = 1 - \left[ \binom{m}{1} - \binom{m}{2} +
\binom{m}{3} - \cdots + (-1)^{m+1}\binom{m}{m} \right] \] Köşeli parantezin
içindeki ifade tam olarak aradığımız net sayım değeridir. Buradan: \[ \text{Net
Sayım} = 1 \] elde edilir. Demek ki bir eleman kaç kümenin kesişiminde olursa
olsun ($m=1, 2, \dots, n$), formüldeki ekleme ve çıkarmalar birbirini götürür ve
o eleman sonuca \textbf{tam 1 kez} eklenmiş olur.

Eleman hiçbir kümede yer almıyorsa ($m=0$), zaten hiçbir terimde sayılmaz,
katkısı 0'dır. İspat tamamlanmıştır.

\subsection{Tümleyen Sayımı (Complementary Counting)}

Çoğu problemde "en az bir koşulu sağlayanlar" yerine, "verilen koşulların
\textbf{hiçbirini} sağlamayanlar" sorulur.

Evrensel kümemiz $U$ olsun. $A_1, A_2, \dots, A_n$ kümeleri de istenmeyen
durumları temsil etsin. Bu durumda hiçbir istenmeyen duruma düşmeyen
elemanların kümesi: \[ |\overline{A_1} \cap \overline{A_2} \cap \dots \cap
\overline{A_n}| = |U| - |A_1 \cup A_2 \cup \dots \cup A_n| \] olarak ifade
edilir. İçerme-dışarma formülünü buraya uygularsak: \[ |U \setminus (A_1 \cup
\dots \cup A_n)| = |U| - \sum |A_i| + \sum |A_i \cap A_j| - \sum |A_i \cap A_j
\cap A_k| + \cdots \] eşitliğini elde ederiz. Bu yaklaşım, özellikle doğrudan
sayımın zor olduğu durumlarda işlemi belirgin biçimde kolaylaştırır.

\begin{definition}[Yaygın Hata Uyarısı] İçerme-dışarma uygularken en sık yapılan
hata terimlerin işaretlerini şaşırmaktır: \begin{itemize} \item \textbf{Birleşim
arıyorsak:} Tekliler +, ikililer -, üçlüler +, dörtlüler - şeklinde ilerler.
\item \textbf{Tümleyen (hiçbirini sağlamayanlar) arıyorsak:} Evrensel küme $|U|$
ile başlarız (+), teklileri çıkarırız (-), ikilileri ekleriz (+), üçlüleri
çıkarırız (-). \end{itemize} Son terimin işaretini ezberlemek yerine, tekli
adımlarda işaretin pozitif mi negatif mi olduğunu takip ederek ilerleyiniz.
\end{definition}

\subsection{Çözümlü Örnekler}

\begin{example} 1'den 1000'e kadar olan tam sayılar arasından 2, 3, 5 veya 7
sayılarından \textbf{hiçbirine bölünmeyen} kaç sayı vardır? \end{example}

\begin{proof}[Çözüm]
Evrensel kümemiz $U = \{1, 2, \dots, 1000\}$ ve $|U| = 1000$'dir. İstenmeyen
durumlar için kümeler tanımlayalım: \begin{itemize} \item $A_1$: 2 ile bölünenler
\item $A_2$: 3 ile bölünenler \item $A_3$: 5 ile bölünenler \item $A_4$: 7 ile
bölünenler \end{itemize} İstenen değer $|\overline{A_1} \cap \overline{A_2} \cap
\overline{A_3} \cap \overline{A_4}|$'tür. Tümleyen formülüne göre: \[ \text{Sonuç}
= 1000 - S_1 + S_2 - S_3 + S_4 \] Burada $S_k$, $k$'lı kesişimlerin eleman
sayıları toplamıdır.

\textbf{1. Adım: Tekli Kümeler ($S_1$)} \begin{align*} |A_1| &= \lfloor 1000 / 2
\rfloor = 500 \\ |A_2| &= \lfloor 1000 / 3 \rfloor = 333 \\ |A_3| &= \lfloor 1000
/ 5 \rfloor = 200 \\ |A_4| &= \lfloor 1000 / 7 \rfloor = 142 \\ S_1 &= 500 + 333
+ 200 + 142 = 1175 \end{align*}

\textbf{2. Adım: İkili Kesişimler ($S_2$)} İkişerli çarpımlara bölünenleri sayarız
($\binom{4}{2} = 6$ terim): \begin{itemize} \item
$|A_1 \cap A_2| = \lfloor 1000 / 6 \rfloor = 166$ \item
$|A_1 \cap A_3| = \lfloor 1000 / 10 \rfloor = 100$ \item
$|A_1 \cap A_4| = \lfloor 1000 / 14 \rfloor = 71$ \item
$|A_2 \cap A_3| = \lfloor 1000 / 15 \rfloor = 66$ \item
$|A_2 \cap A_4| = \lfloor 1000 / 21 \rfloor = 47$ \item
$|A_3 \cap A_4| = \lfloor 1000 / 35 \rfloor = 28$ \end{itemize} \[ S_2 = 166 + 100
+ 71 + 66 + 47 + 28 = 478 \]

\textbf{3. Adım: Üçlü Kesişimler ($S_3$)} Üçerli çarpımlara bölünenleri sayarız
($\binom{4}{3} = 4$ terim): \begin{itemize} \item
$|A_1 \cap A_2 \cap A_3| = \lfloor 1000 / (2 \times 3 \times 5) \rfloor = \lfloor 1000 / 30 \rfloor = 33$
\item
$|A_1 \cap A_2 \cap A_4| = \lfloor 1000 / (2 \times 3 \times 7) \rfloor = \lfloor 1000 / 42 \rfloor = 23$
\item
$|A_1 \cap A_3 \cap A_4| = \lfloor 1000 / (2 \times 5 \times 7) \rfloor = \lfloor 1000 / 70 \rfloor = 14$
\item
$|A_2 \cap A_3 \cap A_4| = \lfloor 1000 / (3 \times 5 \times 7) \rfloor = \lfloor 1000 / 105 \rfloor = 9$
\end{itemize} \[ S_3 = 33 + 23 + 14 + 9 = 79 \]

\textbf{4. Adım: Dörtlü Kesişim ($S_4$)} Dördünün çarpımına
($2 \times 3 \times 5 \times 7 = 210$) bölünenler: \[ S_4 = |A_1 \cap A_2 \cap A_3
\cap A_4| = \lfloor 1000 / 210 \rfloor = 4 \]

\textbf{5. Adım: Birleştirme} \[ \text{Sonuç} = 1000 - 1175 + 478 - 79 + 4 = 228
\] Koşulları sağlayan 228 sayı bulunmaktadır. \end{proof}

\begin{example} $x_1 + x_2 + x_3 + x_4 = 15$ denkleminin, her $i \in \{1, 2, 3, 4\}$
için $0 \le x_i \le 6$ koşulunu sağlayan kaç farklı tam sayı çözümü vardır?
\end{example}

\begin{proof}[Çözüm]
Üst sınır kısıtı ($x_i \le 6$) olmasaydı, problem Bölüm 5'te ele aldığımız
klasik negatif olmayan tam sayılarda ayraç yöntemiyle çözülebilirdi. Koşulu
bozan durumlar, herhangi bir değişkenin $x_i \ge 7$ olmasıdır. Bu istenmeyen
durumları içerme-dışarma ile ayıklayabiliriz.

Evrensel küme $U$, hiçbir üst sınır olmaksızın $x_1 + x_2 + x_3 + x_4 = 15$
denkleminin $x_i \ge 0$ çözümlerinin sayısıdır: \[ |U| = \binom{15 + 4 - 1}{4 - 1}
= \binom{18}{3} = \frac{18 \times 17 \times 16}{6} = 816 \]

İstenmeyen durum kümeleri olarak $A_i$, $x_i \ge 7$ olan çözümler kümesini göstersin
($i = 1, 2, 3, 4$). İstenen çözüm sayısı: \[ \text{Çözüm Sayısı} = |U| - S_1 + S_2
- S_3 + S_4 \]

\textbf{1. Bir değişkenin kısıtı aşması ($S_1$):} Diyelim ki $x_1 \ge 7$.
Değişken değişimi yapalım: $x_1' = x_1 - 7 \ge 0$. Denklem şuna dönüşür: \[ x_1' +
x_2 + x_3 + x_4 = 15 - 7 = 8 \] Bu denklemin negatif olmayan tam sayılarda çözüm
sayısı: \[ \binom{8 + 4 - 1}{4 - 1} = \binom{11}{3} = \frac{11 \times 10
\times 9}{6} = 165 \] Kısıtı aşabilecek $\binom{4}{1} = 4$ farklı değişken
olduğundan: \[ S_1 = 4 \times 165 = 660 \]

\textbf{2. İki değişkenin kısıtı aşması ($S_2$):} Hem $x_1 \ge 7$ hem de $x_2 \ge 7$
olsun. Değişken değişimiyle $x_1' = x_1 - 7$ ve $x_2' = x_2 - 7$ dersek: \[ x_1' +
x_2' + x_3 + x_4 = 15 - 7 - 7 = 1 \] Çözüm sayısı: \[ \binom{1 + 4 - 1}{4 - 1} =
\binom{4}{3} = 4 \] İki değişken $\binom{4}{2} = 6$ farklı şekilde seçilebildiğinden:
\[ S_2 = 6 \times 4 = 24 \]

\textbf{3. Üç veya daha fazla değişkenin kısıtı aşması ($S_3, S_4$):} Üç
değişken birden en az 7 olursa, toplamları en az $7 \times 3 = 21$ olur. Ancak tüm
değişkenlerin toplamı 15 verildiğinden bu durum olanaksızdır. Dolayısıyla
$S_3 = 0$ ve $S_4 = 0$'dır.

\textbf{Sonuç:} \[ \text{Çözüm Sayısı} = |U| - S_1 + S_2 = 816 - 660 + 24 = 180 \]
olarak hesaplanır. \end{proof}

\section{Düzensiz Dizilimler (Derangements) ve Şapka Problemi}

\subsection{Şapka Problemi: Sabit Noktasız Permütasyonlar}

Bölüm 4'te incelediğimiz permütasyon kavramının klasik uygulamalarından biri
vestiyer ya da şapka problemidir: \begin{quote} $n$ kişi bir restorana gider ve
şapkalarını vestiyere teslim eder. Çıkışta vestiyer görevlisi şapkaları rastgele
bir biçimde bu $n$ kişiye dağıtır. \textbf{Hiç kimsenin kendi şapkasını almama}
olasılığı veya bunu sağlayan durum sayısı kaçtır? \end{quote}

Matematiksel olarak bu soru, $\{1, 2, \dots, n\}$ kümesinin hiçbir $i$ elemanı
için $\pi(i) = i$ koşulunu sağlamayan $\pi$ permütasyonlarını bulmaya karşılık
gelir. Bir permütasyonda $\pi(i) = i$ eşitliğini sağlayan elemanlara \textbf{sabit
nokta} (fixed point) denir. Biz sabit noktası bulunmayan permütasyonların
sayısını aramaktayız.

\begin{definition}[Düzensiz Dizilim / Derangement] $\{1, 2, \dots, n\}$ kümesinin
hiçbir elemanını kendi orijinal yerinde bırakmayan permütasyonlarına
\textbf{düzensiz dizilim} (derangement) denir. $n$ elemanlı bir kümenin düzensiz
dizilimlerinin sayısı $!n$ veya $D_n$ simgesiyle gösterilir. \end{definition}

\subsection{Küçük Değerleri Elle İnceleyelim}

Genel formüle geçmeden önce küçük $n$ değerleri için olası permütasyonları
inceleyelim:

\begin{itemize} \item $\mathbf{n = 1}$: Tek bir eleman vardır: $(1)$. Kendi
yerindedir. Sabit noktası olmayan dizilim yoktur: \[ !1 = 0 \] \item $\mathbf{n =
2}$: Olası dizilimler: $(1, 2)$ ve $(2, 1)$. \begin{itemize} \item $(1, 2) \to$ her
ikisi de kendi yerinde. \item $(2, 1) \to 1$, $2$'nin yerinde; $2$, $1$'in yerinde.
İkisi de yerinde değil. \end{itemize} \[ !2 = 1 \] \item $\mathbf{n = 3}$: Toplam
$3! = 6$ permütasyon vardır: \begin{align*} (1, 2, 3) &\to \text{3 sabit nokta}
\\ (1, 3, 2) &\to \text{1 sabit nokta (1)} \\ (3, 2, 1) &\to \text{1 sabit
nokta (2)} \\ (2, 1, 3) &\to \text{1 sabit nokta (3)} \\ (2, 3, 1) &\to
\text{Sabit nokta yok} \quad \checkmark \\ (3, 1, 2) &\to \text{Sabit nokta yok}
\quad \checkmark \end{align*} 6 permütasyondan yalnızca 2 tanesinde hiçbir eleman
kendi yerinde değildir: \[ !3 = 2 \] \item $\mathbf{n = 4}$: Toplam $4! = 24$
permütasyon bulunur. Bunları elediğimizde 9 tanesinin sabit noktası olmadığını
görürüz: \[ !4 = 9 \] \end{itemize} Dizi $0, 1, 2, 9, \dots$ biçiminde ilerler.

\subsection{İçerme-Dışarma ile Genel Formülün Türetilmesi}

Genel $!n$ formülü, içerme-dışarma ilkesinin doğrudan bir sonucudur.

\begin{theorem}[Düzensiz Dizilim Formülü] $n \ge 1$ tam sayısı için düzensiz
dizilimlerin sayısı: \[ !n = n! \left( 1 - \frac{1}{1!} + \frac{1}{2!} -
\frac{1}{3!} + \frac{1}{4!} - \cdots + \frac{(-1)^n}{n!} \right) = n!
\sum_{k=0}^n \frac{(-1)^k}{k!} \] formülü ile hesaplanır. \end{theorem}

\begin{proof} Tüm permütasyonların kümesi evrensel kümemiz olsun: $|U| = n!$.

Bir permütasyonun koşulu ihlal etmesi, en az bir elemanın kendi yerinde kalması
demektir. $i \in \{1, 2, \dots, n\}$ için: \[ A_i = \{ \pi \mid \pi(i) = i \} \]
kümesini tanımlayalım. $A_i$, $i$. elemanın kendi yerinde kaldığı permütasyonların
kümesidir.

Düzensiz dizilimler, bu durumların hiçbirinin gerçekleşmediği
permütasyonlardır: \[ !n = |U \setminus (A_1 \cup A_2 \cup \dots \cup A_n)| \]
Tümleyen içerme-dışarma ilkesini yazalım: \[ !n = |U| - \sum_{i} |A_i| + \sum_{i
< j} |A_i \cap A_j| - \sum_{i < j < k} |A_i \cap A_j \cap A_k| + \cdots + (-1)^n
|A_1 \cap \dots \cap A_n| \] Kesişimlerin büyüklüklerini adım adım hesaplayalım:
\begin{enumerate} \item $|A_i|$: $i$. eleman kendi yerine sabitlenmiştir. Kalan
$n-1$ eleman kendi arasında $(n-1)!$ farklı şekilde sıralanabilir: \[ |A_i| =
(n-1)! \] Bu türden $\binom{n}{1}$ küme olduğundan: \[ \sum_{i} |A_i| =
\binom{n}{1} (n-1)! = \frac{n!}{1!(n-1)!} (n-1)! = \frac{n!}{1!} \] \item
$|A_i \cap A_j|$: Hem $i$ hem de $j$ kendi yerindedir (2 eleman sabit). Kalan $n-2$
eleman $(n-2)!$ biçimde dizilir: \[ |A_i \cap A_j| = (n-2)! \] Bu türden
$\binom{n}{2}$ çift olduğundan: \[ \sum_{i < j} |A_i \cap A_j| = \binom{n}{2}
(n-2)! = \frac{n!}{2!(n-2)!} (n-2)! = \frac{n!}{2!} \] \item Benzer biçimde,
herhangi $k$ elemanın kendi yerinde sabitlendiği durumların sayısı: \[ \sum_{i_1 <
\dots < i_k} |A_{i_1} \cap \dots \cap A_{i_k}| = \binom{n}{k} (n-k)! =
\frac{n!}{k!(n-k)!} (n-k)! = \frac{n!}{k!} \] \end{enumerate} Bu terimleri
formülde yerine yazarsak: \begin{align*} !n &= n! - \frac{n!}{1!} +
\frac{n!}{2!} - \frac{n!}{3!} + \cdots + (-1)^n \frac{n!}{n!} \\ &= n! \left( 1 -
\frac{1}{1!} + \frac{1}{2!} - \frac{1}{3!} + \cdots + \frac{(-1)^n}{n!} \right)
\end{align*} elde edilir. \end{proof}

Formülü küçük değerlerde kontrol edelim: \begin{align*} !1 &= 1!(1 - 1) = 0 \\
!2 &= 2!\left(1 - 1 + \frac{1}{2}\right) = 2 \times \frac{1}{2} = 1 \\ !3 &= 6\left(1
- 1 + \frac{1}{2} - \frac{1}{6}\right) = 6 \times \frac{2}{6} = 2 \\ !4 &= 24\left(1
- 1 + \frac{1}{2} - \frac{1}{6} + \frac{1}{24}\right) = 24 \times \frac{12 - 4 +
1}{24} = 9 \\ !5 &= 120\left(\frac{9}{24} - \frac{1}{120}\right) = 5 \times 9 - 1 =
44 \end{align*} Sonuçlar elle bulduğumuz değerlerle bütünüyle örtüşmektedir.

\subsection{Doğal Bir Sabit: 1/e ile Bağlantı}

Parantez içindeki ifadeye dikkatle bakalım: \[ 1 - \frac{1}{1!} + \frac{1}{2!} -
\frac{1}{3!} + \frac{1}{4!} - \cdots + \frac{(-1)^n}{n!} \] Bu sonlu toplam, $n$
büyüdükçe belirli bir sabite yaklaşır: \[ \frac{1}{e} \approx 0.367879\dots \]
(Bu sabitin neden tam olarak $\frac{1}{e}$ olduğunu merak edenler için hemen
aşağıda kısa bir kutu var.) Bu bağıntı çarpıcı bir gerçeği ortaya koyar: $n$
büyüdükçe rastgele dağıtılan şapkaların \textbf{hiçbirinin} doğru sahibine gitmeme
olasılığı: \[ \frac{!n}{n!} \approx \frac{1}{e} \approx \%36.8 \] değerine hızla
yaklaşır. Kişi sayısı ne kadar büyük olursa olsun, kimsenin kendi şapkasını
alamama olasılığı yaklaşık üçte birdir ($\approx \%36.8$).

\fbox{\parbox{0.92\textwidth}{
\textbf{Meraklısı için: Bu sabit neden tam olarak $1/e$?} \\
$\frac{!n}{n!}$ oranının neden $\frac{1}{e}$'ye yaklaştığını görmek isteyenler için
kısa bir tur: $e^x$ fonksiyonunun Taylor serisi \[ e^x = 1 + \frac{x}{1!} +
\frac{x^2}{2!} + \frac{x^3}{3!} + \cdots \] biçimindedir. Bu seride $x = -1$
alınırsa \[ \frac{1}{e} = e^{-1} = 1 - \frac{1}{1!} + \frac{1}{2!} - \frac{1}{3!}
+ \frac{1}{4!} - \cdots \] elde edilir. İşte düzensiz dizilim formülündeki
parantezli ifade, bu sonsuz toplamın sonlu bir kesitidir; bu yüzden $\frac{!n}{n!}$
oranı, $n$ büyüdükçe $\frac{1}{e}$'ye hızla yaklaşır.}}

Pratik hesaplamalarda $!n$, $\frac{n!}{e}$ değerine en yakın tam sayıdır: \[ !n
= \left\lfloor \frac{n!}{e} + \frac{1}{2} \right\rfloor \quad (n \ge 1) \]

\subsection{Çözümlü Örnekler}

\begin{example} 5 mektup ve bu mektupların gönderileceği 5 farklı adresli zarf
vardır. Mektuplar zarflara rastgele konulduğunda \textbf{tam olarak 2} mektubun
doğru zarfa konulmuş olma sayısı kaçtır? \end{example}

\begin{proof}[Çözüm]
Bu durum iki aşamada incelenebilir: \begin{enumerate} \item Doğru zarfa gidecek
2 mektup seçilir. \item Kalan 3 mektup, hiçbirinin doğru zarfa gitmeyeceği
şekilde dağıtılır. (Eğer kalanlardan biri bile doğru zarfa giderse, doğru zarf
sayısı 2'yi aşacaktır.) \end{enumerate}

Doğru zarfa gidecek 2 mektup $\binom{5}{2}$ farklı şekilde seçilebilir: \[
\binom{5}{2} = \frac{5 \times 4}{2} = 10 \] Kalan 3 mektubun kalan 3 zarfa
tamamen yanlış dağılması ise düzensiz dizilim sayısı olan $!3$ kadardır: \[ !3 = 2
\] Bölüm 3'teki çarpma kuralı gereği: \[ \binom{5}{2} \times !3 = 10 \times 2 = 20
\] farklı durum bulunur. \end{proof}

\begin{example} Bir yuvarlak masa etrafında 4 evli çift (toplam 8 kişi)
oturacaktır. Kadınlar ve erkekler dönüşümlü (bir kadın, bir erkek) oturacaktır.
Kadınlar masaya yerleştikten sonra, hiçbir erkeğin kendi eşinin \textbf{hemen
sağında} oturmamasını istiyoruz. Erkekler bu masaya kaç farklı biçimde
oturabilir? \end{example}

\begin{proof}[Çözüm]
Kadınların yerleri belirlendikten sonra aralarında erkekler için 4 boş sandalye
kalır. Bu sandalyeleri 1, 2, 3, 4 olarak numaralandıralım; öyle ki $i$. sandalye
tam olarak $i$. kadının sağındaki sandalye olsun.

$i$. erkeğin kendi eşinin sağına oturması, $i$. sandalyeye oturması demektir.
Hiçbir erkeğin kendi eşinin sağında oturmaması koşulu ise hiçbir $i$. erkeğin $i$.
sandalyeye oturmaması anlamına gelir.

Bu durum, 4 erkeğin 4 sandalyeye yerleşiminde sabit noktasız bir permütasyon,
yani bir düzensiz dizilim problemidir.

O hâlde aranan oturma düzeni sayısı doğrudan $!4$ ile verilir: \[ !4 = 9 \]
Erkekler, kadınların konumu sabitken bu koşula uygun 9 farklı biçimde
oturabilirler. \end{proof}

\begin{example} $\{1, 2, 3, 4, 5, 6\}$ kümesinin permütasyonları arasından 1'in
ilk sırada olmadığı ($\pi(1) \ne 1$) ve 6'nın son sırada olmadığı ($\pi(6) \ne 6$)
kaç farklı permütasyon vardır? \end{example}

\begin{proof}[Çözüm]
Bu problem tüm elemanların değil, yalnızca iki elemanın konumunu kısıtlamaktadır.
Bu nedenle tam bir düzensiz dizilim formülü yerine iki küme içerme-dışarma
ilkesini uygulamak yeterlidir.

Tüm permütasyonların sayısı $|U| = 6! = 720$'dir.

İstenmeyen durumlar: \begin{itemize} \item $A$: 1'in birinci sırada olması
($\pi(1) = 1$). 1 sabitlendiğinden kalan 5 eleman $5! = 120$ şekilde dizilir:
$|A| = 120$. \item $B$: 6'nın altıncı sırada olması ($\pi(6) = 6$). 6 sabitlendiğinden
kalan 5 eleman $5! = 120$ şekilde dizilir: $|B| = 120$. \item $A \cap B$: Hem
$\pi(1) = 1$ hem de $\pi(6) = 6$ olması. Bu iki eleman sabitlendiğinden kalan 4
eleman $4! = 24$ şekilde dizilir: $|A \cap B| = 24$. \end{itemize} İçerme-dışarma
kuralına göre istenmeyen durumlardan en az birine düşen permütasyon sayısı: \[
|A \cup B| = |A| + |B| - |A \cap B| = 120 + 120 - 24 = 216 \] Her iki ihlalden de
kaçınmak istediğimize göre, evrensel kümeden birleşimi çıkarırız: \[ |\overline{A}
\cap \overline{B}| = |U| - |A \cup B| = 720 - 216 = 504 \] Koşulları sağlayan 504
farklı permütasyon vardır. \end{proof}

\begin{example} İçerme-dışarma mantığını algoritmik olarak nasıl ifade ederiz? $N$
tane kümenin kesişim boyutları verildiğinde genel birleşim boyutunu hesaplayan
mantığı gösteriniz. \end{example}

\begin{proof}[Çözüm] Programlama yarışmalarında $N$ genellikle 15--20 gibi küçük
değerler aldığında, tüm alt kümeleri gezmek için \textbf{bitmask} (bitmask)
tekniği kullanılır. Bir maskede 1 olan bitler seçilen kümeleri, seçilen küme
adedinin tek veya çift olması ise formüldeki + ya da - işaretini belirler:

\begin{lstlisting}
function inclusionExclusion(n, intersectionSize)
    totalUnion = 0
    
    for mask = 1 to (2^n - 1) do
        elementCount = countSetBits(mask)
        intersection = intersectionSize[mask]
        
        if elementCount mod 2 == 1 then
            totalUnion = totalUnion + intersection
        else
            totalUnion = totalUnion - intersection
        end if
    end for
    
    return totalUnion
end function\end{lstlisting}
Bu algoritma $2^n - 1$ alt kümenin her birini içerme-dışarma ilkesinin
belirlediği katsayıyla ($(-1)^{k+1}$) toplama katar.
\end{proof}
